\section{La ruta de Tesla: visión, BEV, Transformer y chips propios}

La sección anterior expuso los problemas de las primeras rutas; esta sección se ocupa de Tesla. Lang Xianpeng (郎咸朋) considera que hacia 2018 la conducción autónoma dio un giro decisivo: Tesla se mantuvo firme en no usar mapas de alta precisión ni LiDAR y se orientó hacia la visión pura, BEV, Transformer y chips de desarrollo propio. Lo esencial no es ninguno de esos términos por separado, sino llevar la información de varias cámaras a una misma representación espacial para que el sistema comprenda de forma unificada el mundo que rodea al vehículo.

\begin{figure}[H]
\centering
\includegraphics[width=0.96\textwidth]{tesla-updimension.png}
\caption{La solución de Tesla mediante elevación de dimensión: resolver problemas del sistema con datos de la flota, visión, chips y entrenamiento a gran escala. Diagrama conceptual propio, elaborado a partir de la conversación de 00:04:30--00:25:06.}
\end{figure}

\begin{knowledgebox}{Lectura de la figura: elevar la dimensión no es acumular más parches}
En la figura, Fleet Data, Vision, Chip, Training y FSD forman una sola línea. La «elevación de dimensión» de Tesla consiste en dejar de corregir cada error de fusión tardía y, en cambio, modificar la forma de representación y los límites del sistema: resolver el problema con un espacio unificado, un modelo unificado, datos de la flota y capacidad de cómputo propia.
\end{knowledgebox}

\subsection{Por qué la visión}

Esta subsección explica la ruta de visión. Elon Musk suele justificar la visión pura con el argumento de que «las personas manejan solo con los ojos», lo cual funciona para la divulgación; en el plano técnico, lo más importante es que las imágenes tienen una alta densidad de información: contienen color, textura, semántica, bordes y geometría implícita. Mediante la paralaje entre varias cámaras, el movimiento entre fotogramas consecutivos y el aprendizaje del modelo, el sistema puede inferir a partir de las imágenes el espacio tridimensional, las velocidades y las tendencias.

\begin{figure}[H]
\centering
\includegraphics[width=0.92\textwidth]{camera-lidar-tradeoff.png}
\caption{La disyuntiva entre Camera y LiDAR: LiDAR proporciona la distancia de forma directa; Camera aporta semántica y escala a bajo costo. Diagrama conceptual propio, elaborado a partir de la conversación de 00:25:06--00:28:46.}
\end{figure}

\begin{knowledgebox}{Lectura de la figura: medir la distancia directamente no equivale a tener más información}
La ventaja de LiDAR es la precisión de la distancia en cada punto; la de Camera es la riqueza de densidad de píxeles y de información semántica. El juicio de Lang Xianpeng es que el LiDAR solo entrega una nube de puntos de alrededor de 128 líneas, mientras que la cámara entrega un mundo a color de varios millones de píxeles. La conducción autónoma no necesita solo distancias, sino también objetos, semántica, intenciones de movimiento y el contexto del entorno.
\end{knowledgebox}

\subsection{Qué resuelven BEV y Transformer}

Esta subsección explica BEV. BEV significa Bird's Eye View, vista de pájaro. El enfoque ingenuo consiste en que cada cámara reconozca primero los objetivos por su cuenta y luego se fusionen los resultados a posteriori; el problema es que las distintas cámaras pueden omitir detecciones, superponerse o tener escalas inconsistentes, y la fusión tardía es propensa a errores. La idea de BEV es extraer primero las características de las múltiples cámaras, unificar la información en una sola representación espacial y luego reconocer de una vez personas, vehículos, carriles y obstáculos. Aquí Transformer es una arquitectura de modelado de propósito general; lo importante es que está al servicio de la comprensión espacial unificada.

\begin{importantbox}{Lo esencial de BEV no es que la vista de pájaro se vea bien, sino que evita los errores de fusión tardía}
La fusión tardía es «cada cámara juzga primero de forma independiente y luego se ensamblan los resultados»; BEV se acerca más a «primero se reúnen las características de todas las cámaras y luego se juzga en un espacio unificado». Esto reduce los conflictos entre resultados de distintas fuentes y mejora la consistencia de los objetos.
\end{importantbox}

\subsection{Fusión temprana y fusión tardía: por qué la idea importa más que el operador}

Esta subsección amplía la explicación de Lang Xianpeng sobre BEV. En los primeros esquemas de varias cámaras, el enfoque ingenuo habitual era que cada cámara ejecutara su propia percepción hacia adelante: la cámara frontal reconocía los carriles y vehículos de enfrente, las cámaras laterales los objetivos a los costados y la cámara trasera los objetivos de atrás; después, el sistema de ingeniería fusionaba esas detecciones en un mismo sistema de coordenadas. El problema es que el resultado de cada cámara puede estar equivocado, incompleto, duplicado o tener una escala inconsistente, y el posprocesamiento debe decidir en qué cámara confiar, si dos mitades de vehículo son el mismo vehículo y cómo empalmar los bordes de las oclusiones. Esto consume una gran cantidad de esfuerzo en «eliminar los errores de la fusión tardía».

La elevación de dimensión de BEV consiste en adelantar la fusión al nivel de las características: primero se extraen características de las imágenes de las múltiples cámaras y luego se infieren objetos, carriles y zonas transitables en un espacio unificado. No se trata simplemente de «unir las imágenes en una imagen grande», sino de que el modelo use al mismo tiempo, en un mismo espacio, la información de varios puntos de vista. El docente enfatiza que Transformer es solo la estructura de red que después se puso al servicio de esta idea; lo verdaderamente decisivo fue haber identificado antes el problema de fondo de la industria de la conducción autónoma: la percepción desde múltiples puntos de vista debe modelarse de forma unificada, no juzgarse localmente en cada vista para luego ensamblarse.

\begin{knowledgebox}{Nota de clase: la lesson de BEV puede trasladarse a muchos sistemas de IA}
Cuando un sistema tiene varias fuentes de información, la fusión tardía suele traer problemas de consistencia. La solución de fondo es hacer una fusión temprana en un espacio de representación adecuado, para que el modelo aproveche todo el contexto de una sola vez. Los modelos multimodales, el estado de las herramientas de un Agent y los VLA de robótica siguen una lógica similar.
\end{knowledgebox}

\begin{figure}[H]
\centering
\includegraphics[width=0.92\textwidth]{av-terms-map.png}
\caption{Mapa de términos clave de la conducción autónoma: la relación entre BEV, Transformer, E2E, VLM, World Model y RL. Diagrama conceptual propio, elaborado a partir de la conversación de 00:20:06--01:26:40.}
\end{figure}

\begin{knowledgebox}{Lectura de la figura: los términos no forman una lista de elementos paralelos}
BEV y Transformer resuelven principalmente la representación y el modelado espaciales; E2E deja al modelo el mapeo de la entrada a la trayectoria; VLM añade comprensión semántica; World Model predice estados futuros; RL permite que el sistema se optimice a partir de la retroalimentación. Al leer esta figura, hay que verla como una cadena que va de la percepción a la decisión y luego al aprendizaje en lazo cerrado.
\end{knowledgebox}

\subsection{Por qué desarrollar chips propios}

Esta subsección trata los chips. El BEV con varias cámaras y el cómputo de modelos grandes requieren mucha capacidad de cómputo en el vehículo. En sus inicios, Tesla pasó de los chips de NVIDIA a su propio chip FSD con esta lógica: solo diseñando en conjunto el algoritmo, los sensores y el chip se obtiene capacidad de cómputo efectiva a un costo relativamente bajo. En la entrevista se menciona que el costo de sensores más chip de Tesla ronda el orden de 1,000 dólares; lo importante no es la cifra exacta, sino mostrar que una ruta de producción en serie debe diseñar de manera conjunta la capacidad de cómputo, el costo y la adaptación del algoritmo.

\begin{knowledgebox}{Términos clave: TOPS, ASIC, capacidad de cómputo efectiva}
\noindent\begin{tabularx}{\textwidth}{p{0.18\textwidth}p{0.36\textwidth}X}
\toprule
Término & Significado & Papel en este episodio \\
\midrule
TOPS & Billones de operaciones por segundo; indicador habitual de la capacidad de cómputo de los chips de IA para vehículos & Un valor alto no garantiza que el modelo real funcione bien. \\
ASIC & Circuito integrado de aplicación específica, hecho a la medida de un algoritmo o carga de trabajo concretos & Tesla puede adaptar mutuamente su algoritmo y su chip propios. \\
Chip de propósito general & Procesador que admite varias empresas y diversos algoritmos & Es flexible, pero no necesariamente óptimo para un algoritmo dado. \\
Capacidad de cómputo efectiva & Capacidad de cómputo que el modelo aprovecha realmente & Se acerca más al resultado de ingeniería que los TOPS nominales. \\
\bottomrule
\end{tabularx}
\end{knowledgebox}

\subsection{Resumen de la sección}

La esencia de la ruta de Tesla es llevar la conducción autónoma de la certeza basada en sensores costosos y mapas hacia la visión, la representación espacial unificada, los datos de la flota, los chips propios y el entrenamiento a gran escala. Fue la primera vez que la conducción autónoma se orientó explícitamente hacia un «sistema basado en la capacidad del modelo».
